Fix integers $r,k\ge 2$.
Let $A\subseteq\mathbb{R}^2$ be a set of $n$ points with the property that no line contains $k$ points of $A$.
A line determined by two points of $A$ is called \emph{ordinary} if it contains exactly two points of $A$.
Let $L_\mathrm{ord}(A)$ denote the number of ordinary lines determined by $A$.

Define $f_{r,k}(n)$ to be the smallest integer with the property:

If $|A|=n$, $A$ has no $k$ collinear points, and $L_\mathrm{ord}(A)\ge f_{r,k}(n)$, then there exists a subset $A'\subseteq A$ with $|A'|=r$ such that every line determined by two points of $A'$ is ordinary (in $A$).

Question: estimate $f_{r,k}(n)$. In particular, is $f_{r,k}(n)=o(n^2)$? Even $f_{r,k}(n)=O(n)$?
